\documentclass[10pt, romanian]{article}
\usepackage[romanian]{babel}
\usepackage{natbib}
\usepackage{charter}
\usepackage{fullpage}
\usepackage[colorlinks=false]{hyperref}
\usepackage{ifthen}
\usepackage[resetlabels]{multibib}
%\usepackage[openbib,ManyBibs]{currvita}
\usepackage[ManyBibs]{currvita}
%\usepackage{currvita}
\usepackage[utf8x]{inputenc}
\usepackage{amssymb}

\newboolean{doResearchInterests}
\setboolean{doResearchInterests}{true}
%\setboolean{doResearchInterests}{false}
\newcites{thesis}{A. Teza de doctorat}
\newcites{book}{B. Cărți și capitole în cărți}
\newcites{isij}{C. Lucrări indexate ISI}
\newcites{isic}{C. Lucrări indexate ISI}
\newcites{neindexate}{D. Lucrări publicate în reviste și volume de conferințe cu referenți (neindexate)}

\usepackage{ucs} % for unicode characters (just for \Rosu and \Serbanuta)
\newcommand{\sh}{\unichar{0537}}
\newcommand{\Sh}{\unichar{0536}}
\newcommand{\tz}{\unichar{0539}}
\newcommand{\Tz}{\unichar{0538}}
\PrerenderUnicode{\sh}
\PrerenderUnicode{\Sh}
\PrerenderUnicode{\tz}
\PrerenderUnicode{\Tz}
\newcommand{\Serbanuta}{{\Sh}er\-b{\u a}\-nu\-{\tz}{\u a}}
\newcommand{\Rosu}{Ro\-{\sh}u}
\newcommand{\Stefanescu}{{\Sh}te\-f{\u a}\-nes\-cu}


% Better for lists with 1-2 items and short descriptions
\newenvironment{sublist}{%
    \begin{list}{}{%
        \setlength{\itemsep}{0em}\setlength{\parsep}{0em}%
        \setlength{\topsep}{0em}\setlength{\parskip}{0em}%
    }%
}%
{ \end{list} }

% Better for lists with more than 2 items and/or long descriptions
\newenvironment{subbulletlist}{%
    \begin{list}{\labelitemii}{%
        \setlength{\topsep}{\itemsep}\setlength{\parskip}{\parsep}%
    }%
}%
{ \end{list} }

\begin{document}

%\bibliographystyle{unsrt}
%\bibliographystylebook{unsrt}
%\bibliographystyleconference{unsrt}
%\bibliographystyleworkshop{unsrt}
%\bibliographystylejournal{unsrt}
%\bibliographystyletechreport{unsrt}

\newlength{\oldcvlabelwidth}
\setlength{\oldcvlabelwidth}{\cvlabelwidth}

\renewcommand*{\cvbibname}{}

\begin{cv}{Traian Florin Șerbănuță\\{\large \itshape Curriculum Vit\ae}}

\begin{cvlist}{Date Personale}
   \setlength{\itemsep}{-1ex}
    \item[Date de contact] Adresa: Facultatea de Matematică și Informatică, Universitatea din București
    \\Str. Academiei nr.14, sector 1, C.P. 010014, Bucuresti, Romania
    \item Tel: (4-021) 314 3508, Fax: (4-021) 315 6990
    \item \href{mailto:traian.serbanuta@unibuc.ro}{traian.serbanuta@unibuc.ro}
    \item \href{http://unibuc.ro/~tserbanuta}
    {http://unibuc.ro/~tserbanuta}
       \setlength{\itemsep}{0ex}
\item[Cetățenie] română
\item[Stare Civilă] Căsătorit, trei copii
\end{cvlist}

\begin{cvlist}{Educație}
   \setlength{\itemsep}{-1ex}
    \item[2010] \emph{Doctor în Informatică},
    Universitatea din Illinois, Urbana-Champaign
    %\begin{sublist}
        \item Teza: \textit{A Rewriting Approach to Concurrent Programming Language Design and Semantics}
        \item Coordonator: Grigore Roșu
        \item Comisia de doctorat: Thomas Ball, Darko Marinov, Jos\'e Meseguer,  Madhusudan Parthasarathy
  %  \end{sublist}
   \setlength{\itemsep}{0ex}
    \item[2004] \emph{Master în Informatică},
    Universitatea din București
   \setlength{\itemsep}{-1ex}
%    \begin{sublist}
        \item Teza: \textit{Concepte instituționale în logica de ordinul I, teoria specificațiilor parametrizate și programarea logică}
        \item Coordonatori: Răzvan Diaconescu și Virgil Emil Căzănescu
%    \end{sublist}
   \setlength{\itemsep}{0ex}
    \item[2002] \emph{Licențiat în Informatică},
    Universitatea din București
   \setlength{\itemsep}{-1ex}
%    \begin{sublist}
        \item Teza: \textit{Ascunderea informației în text folosind gramatici de tip LR(k)}
        \item Coordonator: Adrian Atanasiu
 %   \end{sublist}
\end{cvlist}


\ifthenelse{\boolean{doResearchInterests}} {
\setlength{\cvlabelwidth}{1em}
\begin{cvlist}{Interese de cercetare}
\item Logici formale și aplicațiile lor în informatică
\item Specificații algebrice și metode formale pentru ingineria sofware
\item Algoritmi, programare funcțională, programare logică, rescriere
\end{cvlist}
\setlength{\cvlabelwidth}{\oldcvlabelwidth}
} {
}

\begin{cvlist}{Experiență profesională}
    \item[2013--prezent] Conferențiar universitar la {\em Universitatea din București}, Facultatea de Matematică și Informatică, Departamentul de Informatică
    \item[2013--prezent] Consultant/Cercetător pentru Runtime Verification, Inc.
    \item[2011--2013] Cercetător știintific post-doctoral la {\em Universitatea Alexandru Ioan Cuza Iași}\\
    Laboratorul de Metode Formale in Ingineria Software (FMSE)
    \item[2011--2012] Colaborator de cercetare post-doctoral la {\em Universitatea din Illinois, Urbana-Champaign}\\
    Information Trust Institute (ITI), Laboratorul de sisteme formale (FSL)
    \item[2004--2010] Asistent de cercetare la {\em Universitatea din Illinois,
    Urbana-Champaign}\\
    Laboratorul de sisteme formale (FSL)
    \item[2003--2009] Preparator la {\em Universitatea din București}\\
     Catedra de Fundamentele informaticii
     \item[2007] Practică de vară la {\em Google New York}\\
     Co-autor (împreună cu Bogdan Căpriță) al unei aplicații de patent referitor la o metodă de analizare a traficului web
     \item[2005] Practică de vară la {\em Microsoft Research Redmond}\\
     Grupul de Testare, Verificare și Măsurare (TVM)
     \item[2000--2001] Programator la {\em Popnet-Agentscape Romania}\\
     Echipa de procesare a limbajului natural
\end{cvlist}

%\ifthenelse{\boolean{doResearchInterests}} {
%\pagebreak
%}{}

%\setlength{\bibindent}{2.5em}
% Do bibaddtoleftmargin instead of biblabelsep if you're
% not using manybib or openbib
%\renewcommand*{\bibaddtoleftmargin}{1.5em}
%\setlength{\bibhang}{3.5em}
% To update, before doing latexmk, run:
% bibtex journal
% bibtex refpaper
% bibtex unrefpaper
% bibtex techreport
%\newpage
  \begin{cvlist}{Bibliometric data}
  \item[Publications in ISI Journals] 11
  \item[Publications in ISI Proceedings] 19
  \item[Hirsch index] 12 --- ISI Web of Science, 14 --- Scopus, 19 --- Google Scholar
  \item[Citations] 298 -- ISI Web of Science, 719 --- Scopus, 1436 --- Google Scholar
  \end{cvlist}



\setlength{\cvlabelwidth}{\oldcvlabelwidth}

\begin{cvlist}{Prezentări}
\item[CRM9'19] \textit{Operational Semantics and Program Verification Using Many-Sorted Hybrid 
               Modal Logic}
\item[TABLEAUX'19] \textit{Operational Semantics and Program Verification Using Many-Sorted Hybrid
               Modal Logic}
\item[Diaspora și prietenii săi'16] \textit{Maximally Parallel Contextual String Rewriting} la invitația organizatorilor
\item[WRLA'16] \textit{Maximally Parallel Contextual String Rewriting}
\item[SYNASC'12] Susținerea unui tutorial despre K (împreună cu Dorel Lucanu) la invitația organizatorilor
\item[RV'12]  \textit{Maximal Causal Models for Sequentially Consistent Systems}
\item[ICGT'12] \textit{A Truly Concurrent Semantics for the K Framework Based on
               Graph Transformations}
               \\
               De asemeni, susținerea unei prezentari depre K în cadrul sesiunii supriză PechaKucha, la invitația organizatorilor.
\item[K'11] \textit{A Concurrent Semantics for the K Framework}
\\ \textit{From Language Definitions to (Runtime) Analysis Tools}
\\ \textit{The K-framework Tool Chain---towards version 3.0: lessons learned and new perspectives}
\item[WRLA'10] \textit{K-Maude: A Rewriting Based Tool for Semantics of Programming Languages}---\\prezentare și demonstrație
\item[ICSE'08] \textit{JPredictor: A Predictive Runtime Analysis Tool for Java}
\item[WMC'08] \textit{Defining and Executing P-systems with Structured Data in K}
\item[SOS'07] \textit{A Rewriting Logic Approach to Operational Semantics---Extended Abstract}
\item[WRLA'06] Co-organizator și co-prezentator al \textit{Competiției de motoare de rescriere}
\item[FOSSACS'06] \textit{A Semantic Approach to Interpolation}
\item[MSR Redmond'05] \textit{Interleaving Boolean and theory reasoning in a SAT-based theorem prover}---\\prezentare de încheiere a practicii de vară
\end{cvlist}

%\newpage
%\setlength{\cvlabelwidth}{1em}
%
%\begin{cvlist}{Premii}
%\item[2000]%[8/2004--prezent]
%    \emph{Locul 15 (ca membru al echipei Universității București) la faza finală a concursului internațional de programare ACM}
%\item[1999]%[8/2001--8/2004]
%    \emph{Locul întâi (ca membru al echipei Universității București) la faza pe Europa de sud-est al concursului internațional de programare ACM}
%%\item[1998]
%%    \emph{Premiul III la Olimpiada națională de matematică}
%\end{cvlist}
%

\begin{cvlist}{Recenzent}
        \item \textit{Journal of Theoretical Computer Science}, \textit{International Journal of Foundations of Computer Science}, \textit{Science of Computer Programming}, \textit{Journal of Computer and System Sciences}, \textit{Fundamenta Informatic\ae}.
        \item  \input reviewer.tex
\end{cvlist}

%\setlength{\cvlabelwidth}{\oldcvlabelwidth}
%\begin{cvlist}{Extracurricular}
%    \item[2008--2011] Ajutor comunitar la
%    University of Illinois Family and Graduate Housing
%   \setlength{\itemsep}{-1ex}
%%    \begin{subbulletlist}
%    \item Atribuții: Propunerea, organizarea și asistarea programelor pentru  membrii comunității.
% %   \end{subbulletlist}
%   \setlength{\itemsep}{0ex}
%    \item[2010--2012] Membru al consiliului parohial pentru
%    biserica ortodoxă Trei Ierarhi din Champaign, IL
%\end{cvlist}

\setlength{\oldcvlabelwidth}{\cvlabelwidth}
\setlength{\cvlabelwidth}{1em}
\setlength{\cvlabelwidth}{\oldcvlabelwidth}

%\begin{cvlist}{References}
%\item Available upon request
%\item See attached sheets
%\end{cvlist}




\bibliographystyle{plain}
\bibliographystylethesis{unsrt}
\bibliographystylebook{unsrt}
\bibliographystyleisij{unsrt}
\bibliographystyleisic{unsrt}
\bibliographystyleneindexate{unsrt}


\renewcommand*{\cvbibname}{}

  \begin{cvlist}{Lista lucrărilor științifice}
  \item[A. Teza de doctorat]
  \nocitethesis{serbanuta-2010-thesis}
  \bibliographythesis{vitae}
  \item[B. Cărți]
  \input books.tex
  \bibliographybook{vitae}
  \item[C. Lucrări indexate ISI]

  \item[\hfill În reviste]
    \input isi-journal.tex
    \bibliographyisij{vitae}
  \item[\hfill În conferințe]
    \input isi-conference.tex
    \bibliographyisic{vitae}
  \end{cvlist}

\end{cv}
\end{document}


